% RENUMBERED 2026-11-16 (user: "rename the files also"): methodology files follow the reading order.
%   Old -> new: 05 features -> 03 (05a/05b/05c -> 03a/03b/03c), 06 feature validity -> 04,
%   03 statistical model (before that 03_prediction_model) -> 05, 04 model size -> 06; 07 unchanged.
%   File names in the comments below were updated; bare numbers in comments dated before this entry
%   ("Methods 03", "05a header", "05d") are the OLD numbers.
% Methodology, section 5: the statistical model, the test of the geometric features, the secondary
% performance measures and their cross-validation.
% 2026-11-16 (user: "write the prediction_model.tex chapter in methodology. Rename the file ... write what
%   you think is the best. Pull a lot from the technical background. refer back to the chapter like aida
%   does ... warm tone like aida"): file renamed 03_prediction_model.tex -> 03_statistical_model.tex (-> 05_statistical_model.tex later the same day) and
%   the section written. The earlier notes and dated entries are in git (this file before the rename);
%   what is still open from them is listed below.
%   2026-11-16 (user: "move it up"): this section now comes after the geometric features (03_features) and
%   their measurement properties (04_feature_validity) in content/methods.tex, so the opener chains to the features and the
%   feature equations (eq:method-tl, eq:method-angle, eq:method-asymmetry) are back-references.
%   sec:method-features points forward to this section for the proximal LAD and LCx rule.
%   2026-11-16 (user: "yes", remove the model size section): the section "Model Size and Feature Selection"
%   (06_model_size.tex, labels sec:method-size, sec:method-size-budget, sec:method-selection) was retired to
%   outdated/thesis_retired/methodology/retired/06_model_size_2026-11-16.tex. Its one remaining fact, the
%   number of fixed effects, is the last paragraph of sec:method-model-spec: 9 segment intercepts + 6 risk
%   factors + 3 (log follow-up, dominance, log length) = 18 in the base, + 3 features = 21. No threshold
%   and no citation (root CLAUDE.md, Riley criteria dropped). CHECK: redo the count if the red note on
%   follow-up, dominance and length removes a term. Counts (preliminary: 1,076 participants, 452 with
%   incident plaque, 970 segment events) belong to sec:results-cohort, which is not written.
%   CAREFUL: "fixed before the geometric features were related to the outcome" must hold for the
%   clean cohort (open decision 12).
%   STRUCTURE: opener (unit = segment, cluster = participant; the two cloglog sentences of 2026-11-14 and
%   2026-11-16 kept word for word; pointer to app:link added, owed since 2026-11-16), then
%   sec:method-model-spec, sec:method-estimation, sec:method-test (main result), sec:method-added-value
%   (secondary), sec:method-cv, sec:method-sensitivity.
%   THEORY IS POINTED BACK TO, NOT RESTATED: sec:tech-model (GLMM), sec:tech-segment-glmm (link, hazard
%   ratio, per SD, exposure), sec:tech-random-intercept (frailty, MHR, conditional reading),
%   sec:tech-glmm-likelihood (likelihood, eq:tech-lr), (sec:tech-aggregation moved here 2026-11-18 as eq:method-agg-conditional),
%   app:link, app:glmm-fitting, app:glmm-variance (eq:tech-mhr, second random effect), app:glmm-risk
%   (eq:tech-agg-marginal). The appendix pointers this section owed are now all in the text.
%   MOVED OUT: the subsection "Recalibration of SCORE2" (label sec:method-recalibration, eq:method-recal)
%   went to Appendix D (appendix/score2_recalibration.tex), since the recalibrated score is only a
%   sensitivity analysis; sec:method-sensitivity keeps one sentence and points there.
%   VOICE: Aida's Chapter 5 measured (105 sentences, median 24 words; chained openers, reason first,
%   "As described in Section ..."). No "prognostic factor", "block" (the Technical Background says
%   "set"), "budget", causal or risk-model wording.
%
% 2026-11-18 (user: "rewrite the cloglog reason in methods 03"): the two cloglog sentences of 2026-11-14 and
%   2026-11-16 (kept word for word until now) and the appendix pointer replaced. The reason is now a choice of
%   reading, not a requirement: logit can be fitted to the same outcomes, take log follow-up as a covariate, and
%   give a participant risk as 1 - prod(1 - p_ij). Order: grouped time first (the more basic reason), then the
%   segment hazards adding to one participant hazard. NEW: a participant-level covariate multiplies H_i by
%   exp(beta), given b_i, from eq:method-agg-conditional (uncited, follows from the equation; it does not hold
%   exactly for the risk averaged over b_i, so the sentence says "given the participant intercept").
%   Logit sentence rests on app:link-logit (proportional odds; "A product of logistic terms is not logistic").
%   CUT here: \cite[pp.~58--59]{prentice1978grouped}. Not opened, no PDF in papers/; the sentence now points
%   to sec:tech-segment-glmm, whose derivation rests on Dobson 10.2 and 10.2.2. Prentice is still cited in
%   Appendix B (pp. 58-59 and p. 66): CHECK there.
%   Left to the Discussion: proportional hazards and proportional odds cannot be tested with one follow-up scan.
%   Not used as a reason: SCORE2 is recalibrated on the cloglog scale (relevant only to the sensitivity
%   analysis, app:score2-recal).
% WRITTEN AS THE DECIDED DESIGN (user, 2026-11-16), NOT AS THE SCRIPT STANDS. The model script is
%   CGPS_data/progression_who_pooled.py, not in CGPS_data/thesis_final/. The engine matches the text:
%   cloglog GLMM, participant intercept, 30 Gauss-Hermite nodes, refit with doubled nodes (--check,
%   NODE_TOL beta 1e-3, sigma 1e-3), LR test on the full-data fits, Wald intervals (wald_se),
%   participant risk by integration over b (person_risk), 10 folds x 50 repeats by participant with
%   seeds 0..49, constants learned on the training participants, mean and 2.5/97.5 percentiles.
% CHECK, each a point where the script differs from this text; align the script or the text:
%   - Cohort and outcome: the script runs on the Medis trees with progression (higher grade at
%     follow-up) as the outcome and a baseline-grade term. Text: plaque-free cohort, incident plaque.
%   - Base: the script defaults to standardised log SCORE2 (--person score2); the six refitted factors
%     are --person factors. Text: the six factors.
%   - (2026-11-18, resolved) Curvature and angle now come from thesis_final (curvature_feature.
%     artery_curvature, bifurcation_angle.lm_angle) in the script's own build(); mean_T5, the turning
%     variants and progression_angle are gone (files deleted). G02 = curvature + angle + asymmetry.
%   - Daughter asymmetry: |a| on the narrower daughter's row (daughter_ratio.narrower_rows), divided by
%     the training SD, not centred; complete cases (persons without it leave both models, text in
%     sec:method-estimation). The column stays empty until tracer contours can be read, and the script
%     then leaves the term out with a note.
%   - Quadrature: the script uses ordinary Gauss-Hermite nodes (hermgauss, centred at 0), while
%     app:glmm-fitting describes the adaptive form. The text says only "Gauss-Hermite quadrature".
%     Change the appendix sentence or the script. The participant risk uses 40 nodes (not stated).
%   - The script adds a ridge term of 1e-6 on the coefficients (fit_glmm, l2); not stated.
%   - The script also checks the log-likelihood (0.05) and the participant risk (1e-6); the text names
%     only the coefficient and standard-deviation tolerance.
%   - Missing curvature or angle: the script sets the standardised value to 0; the plan (CGPS_statistics-
%     CLAUDE.md, open decision 9) is multiple imputation for the test and z = 0 in cross-validation.
%     Red note in the text.
%   - Not in the script: MHR, back-to-back histograms, curvature x artery term, the mild-or-worse
%     outcome, the diabetes/BMI/height model, the recalibrated score on the cloglog scale (the script
%     enters log SCORE2; score2_honesty.py uses cloglog(SCORE2) at person level).
%   - In the script but not in the text, on purpose: Youden confusion matrices (no clinical threshold
%     exists), per-segment AUC and ROC panels, the ten turning variants (exploratory).
%   - Calibration slope and intercept: the script fits them on the logit of the predicted risk
%     (metrics, cal_intercept with the slope fixed at 1). The text does not name the scale.
% OPEN, red \clarify in the text: whether log follow-up, dominance and log length stay in the base
%   (root CLAUDE.md, Open); the significance level; the rule for missing geometry; the artery-level
%   random effect (fit it, check residual correlation, or state the limitation; if none is fitted or
%   discussed, cut the sentence in sec:tech-glmm-likelihood and the paragraph sec:tech-nested);
%   sampling intervals for the differences (bootstrap of the whole procedure, with folds assigned by
%   original participant, open decision 1).
% OPEN, not in the text: absent left main and ramus/trifurcation (open decision 5); the left main
%   features on the proximal rows make the test partly a test of which segment develops plaque
%   (decision 5), to be said in the Discussion or handled by a main effect plus a proximal contrast;
%   participants with an ungraded segment among the nine at follow-up (decision 7).
% CITATIONS. dobson2002glm Section 10.2.3 opened 2026-11-16 (PDF in the repository root): Weibull
%   model, hazard lambda phi y^(lambda-1) (Eq. 10.12), "with suitable values of lambda it can increase
%   or decrease with increasing survival time", exponential (constant hazard) at lambda = 1,
%   log H = log phi + lambda log y (Eq. 10.13). SETTING: Dobson states this for survival times of one
%   sample; here gamma is read across participants with different follow-up times, which is our step.
%   score2_2021: only for the six factors being those of SCORE2 (as in Data 02).
%   prentice1978grouped pp. 58-59: CHECK, still not opened (no PDF in the repository).
% CITATION WANTED: the C-statistic, calibration slope and intercept, Brier score and log loss (none in
%   refs.bib; the Technical Background has no added-value section, so they are defined here in one
%   clause each); the back-to-back histograms (Harrell, fharrell.com/post/addvalue, or Regression
%   Modeling Strategies 2015; not opened, so not cited); cross-validation and optimism.
% CAREFUL: "prespecified" must stay true: no geometry is related to incident plaque in the clean
%   cohort before the pipeline is frozen (open decision 12).
% 2026-11-18 (user): angle x asymmetry term DROPPED (Garcha p. 7: the angle's effect changes sign with
%   which daughter is wider, so one product term does not test it). sec:method-sensitivity now has one
%   secondary analysis (the descriptive where reading); "First ... Subsequently" gone.
% 2026-11-18 (user, plan approved): asymmetry enters only the narrower daughter's proximal segment as
%   |a_i| (sec:method-ratio). eq:method-lp: m_j(beta_theta theta_i + beta_a a_i) -> m_j beta_theta
%   theta_i + beta_a a_ij; a_ij defined after m_j. Standardisation sentence split: risk factors,
%   curvature, angle centred and scaled; asymmetry only scaled, so 0 keeps meaning equal daughters or the
%   wider one (centring would put an effect on the wider daughter). Table form "absolute value, per SD";
%   level stays "participant" (measured once per participant, like the angle). Fixed-effect count
%   unchanged (21).
% Histograms: written on the cross-validated risks, like the other secondary measures. Harrell's own
%   example uses in-sample risks; change the sentence if that is preferred.
% 2026-11-18 (user: "I will have two models in my thesis a glm for my features that are person level and a
%   glmm for curvature feature"; answers: outcome of the GLM = plaque in the proximal LAD or proximal LCx;
%   cloglog + the six risk factors; two prespecified LR tests; docs + thesis text). Section rewritten for two
%   models, each with a base and an extended fit:
%   - Curvature model: the segment GLMM of before, with curvature as the only feature (eq:method-lp). Base 18
%     fixed effects, extended 19. Test q = 1.
%   - Bifurcation model: a participant-level cloglog GLM (sec:tech-glmi), outcome plaque in the proximal LAD
%     or proximal LCx (defined in sec:method-population), angle and |a| (eq:method-lp-bif). Base: intercept,
%     six factors, log follow-up, dominance = 9; extended 11. Test q = 2. No segment length (a segment term);
%     whether the summed length of the two proximal segments should enter is in the \clarify.
%   - m_j and a_ij (narrower-row coding) are gone from the text; |a_i| per SD, not centred, stays.
%   - Complete cases for the asymmetry now apply to the bifurcation model only.
%   - "This is the only test in the thesis" -> two tests; multiplicity added to the significance \clarify.
%   - The added-value paragraph and the dilution argument (eq:method-agg-conditional) belong to the curvature model;
%     the bifurcation model gives its participant risk directly.
% NOT IN THE CODE (CHECK): progression_who_pooled.py fits only the segment GLMM, with angle and asymmetry on the
%   proximal rows. No participant-level GLM for the bifurcation features exists yet, so sec:method-models
%   describes the decided design, not the script. Write the GLM (or a --model bifurcation option) before
%   this text is final.
% Uncited, from the equations: under the cloglog link, the risk of plaque in the proximal LAD or LCx is
%   1 - exp(-(H_pLAD + H_pLCx)), so the bifurcation model keeps the same link as the summed hazards of the
%   two segments. Not stated in the text.
\section{Statistical Model and Evaluation}
\label{sec:method-models}

Following the computation of the geometric features, the features were related to the outcome in two models, according to the level at which they were measured.
Curvature was measured in each of the nine segments of a participant, each with its own outcome.
The nine segments of a participant share the same risk factors and the same unmeasured traits, so their outcomes are correlated.
Curvature was therefore analysed in the generalised linear mixed model described in Section~\ref{sec:tech-model}, with the segment as the unit and the participant as the cluster, referred to as the curvature model.
The bifurcation angle and the daughter asymmetry were measured once per participant and related to plaque in the proximal LAD or the proximal LCx (Section~\ref{sec:method-population}).
This gives one outcome per participant, so they were analysed in the generalised linear model of Section~\ref{sec:tech-glmi}, referred to as the bifurcation model.

Plaque was observed only at the follow-up scan, so the time at which it formed was unknown.
To describe such an outcome through its hazard, the outcomes of both models were modelled with a complementary log-log link (Section~\ref{sec:tech-segment-glmm}).
Under proportional hazards, each coefficient is then a hazard ratio that holds at every time in the follow-up, so its reading does not depend on the length of the follow-up, which differed between participants.
In the curvature model, the same link makes the hazards of the nine segments of a participant add to one participant hazard (Section~\ref{sec:method-added-value}).
Given the participant intercept, a covariate with one value per participant multiplies the participant hazard by the same factor as each segment hazard, so its hazard ratio holds for the participant as well as for each segment.
Under the logit link, the odds ratio is constant over time only under proportional odds, and the odds of the segments do not add (Appendix~\ref{app:link}).

The remainder of this section describes the two models, how they were fitted, the tests of the geometric features, and the secondary estimates of how much the features add to prediction.

\subsection{Base and Extended Models}
\label{sec:method-model-spec}

Each of the two models was fitted twice, as a base model and as an extended model.
The base model held the established risk factors, and the extended model added the geometric features of that model.
For segment $j$ of participant $i$, the linear predictor of the extended curvature model was
\begin{equation}
  \eta_{ij} = \alpha_j + \sum_{k=1}^{6} \beta_k r_{ik} + \gamma \log E_i + \beta_d d_i + \beta_L \log L_{ij} + \beta_T \tilde{T}_{ij} ,
  \label{eq:method-lp}
\end{equation}
where $\alpha_j$ is the intercept of segment $j$, $r_{ik}$ are the six risk factors, $E_i$ is the follow-up time in years, $d_i$ indicates right dominance and $L_{ij}$ is the length of the segment in mm.
$T_{ij}$ is the average curvature of the segment (Eq.~\eqref{eq:method-tl}), and the tilde marks a standardised feature.
The base curvature model is Eq.~\eqref{eq:method-lp} without the curvature term, and both curvature models held the participant intercept $b_i$ of Eq.~\eqref{eq:tech-glmm}.

For participant $i$, the linear predictor of the extended bifurcation model was
\begin{equation}
  \eta_{i} = \alpha + \sum_{k=1}^{6} \beta_k r_{ik} + \gamma \log E_i + \beta_d d_i + \beta_\theta \tilde{\theta}_i + \beta_a \tilde{a}_i ,
  \label{eq:method-lp-bif}
\end{equation}
where $\alpha$ is the intercept, $\theta_i$ is the bifurcation angle (Eq.~\eqref{eq:method-angle}) and $a_i$ is the absolute daughter asymmetry $|a|$ (Eq.~\eqref{eq:method-asymmetry}).
The base bifurcation model is Eq.~\eqref{eq:method-lp-bif} without the two geometric terms.

The six risk factors were age, sex, current smoking, systolic blood pressure, total cholesterol and HDL cholesterol, which are the risk factors of SCORE2 \cite{score2_2021}.
This set was chosen since it was fixed in advance by an external source.
Each factor entered both models as its own linear term, so its weight was estimated for incident coronary plaque and not taken from the score.
In the curvature model, an intercept per segment gave each of the nine segments its own baseline hazard, so differences in how often plaque develops along the tree were not attributed to the covariates.

The time between the two scans differed between participants, so the follow-up time entered as the exposure of Section~\ref{sec:tech-segment-glmm}, with an estimated coefficient $\gamma$.
A value of $\gamma = 1$ corresponds to a constant hazard over the follow-up.
A value above 1 corresponds to a hazard that rises with time and a value below 1 to one that falls, as for the shape parameter of a Weibull model \cite[Section~10.2.3]{dobson2002glm}.
Finally, both base models held right dominance, since the sizes of the RCA and the LCx depend on it (Section~\ref{sec:dominance}).
The base curvature model also held the logarithm of the segment length, since a longer segment offers more wall on which plaque can be detected.
\clarify{whether the follow-up time, dominance and segment length stay in the base models is not yet decided, and whether the bifurcation model should hold the length of the proximal LAD and LCx.}

In the extended models, each geometric feature entered as a single linear term.
To put their hazard ratios on a common scale (Section~\ref{sec:tech-segment-glmm}), the continuous risk factors, curvature and the bifurcation angle were standardised to a mean of zero and a standard deviation of one.
Curvature was standardised within each of the nine segments.
Its hazard ratio therefore compares a segment with the same segment in other participants, and one hazard ratio per standard deviation is assumed for all nine.
The absolute daughter asymmetry was only divided by its standard deviation over the participants, so that zero still marks equal daughters.
The covariates of the models are summarised in Table~\ref{tab:method-covariates}.

\begin{table}[htbp]
  \centering
  \caption{Summary of the covariates of the curvature and the bifurcation model, including the level at which each is measured, the form in which it entered and the fits that held it.}
  \label{tab:method-covariates}
  \begin{tabular}{llll}
    \toprule
    Covariate & Level & Form & Model \\
    \midrule
    Segment & segment & nine intercepts & curvature, both fits \\
    Age, systolic blood pressure & participant & linear, per SD & all fits \\
    Total and HDL cholesterol & participant & linear, per SD & all fits \\
    Sex, current smoking & participant & yes or no & all fits \\
    Follow-up time & participant & logarithm & all fits \\
    Right dominance & participant & yes or no & all fits \\
    Segment length & segment & logarithm & curvature, both fits \\
    Curvature & segment & linear, per SD within segment & curvature, extended \\
    Bifurcation angle & participant & linear, per SD & bifurcation, extended \\
    Daughter asymmetry & participant & absolute value, per SD & bifurcation, extended \\
    \bottomrule
  \end{tabular}
\end{table}

The level of a covariate decides how its hazard ratio is read in the curvature model (Section~\ref{sec:tech-random-intercept}).
Curvature differs between the segments of a participant, whereas the risk factors have one value per participant.

In total, the base curvature model held 18 fixed effects: the nine segment intercepts, the six risk factors, and the follow-up time, dominance and segment length.
The extended curvature model added curvature, which gives 19.
The base bifurcation model held 9 coefficients: the intercept, the six risk factors, the follow-up time and dominance.
The extended bifurcation model added the bifurcation angle and the daughter asymmetry, which gives 11.
All of them were fixed before the geometric features were related to the outcome, and the numbers of participants and segments with incident coronary plaque behind them are reported in Section~\ref{sec:results-cohort}.

\subsection{Estimation}
\label{sec:method-estimation}

All fits were by maximum likelihood on the participants of the study cohort (Section~\ref{sec:method-population}).
The base and the extended fit of a model always used the same segments or participants, so that their likelihoods could be compared.
Participants without a daughter asymmetry were left out of both bifurcation fits, since a value of zero would have marked them as having equal daughters.
In the curvature model, the likelihood is averaged over the unobserved participant intercepts, as described in Section~\ref{sec:tech-glmm-likelihood}.
This average was computed by Gauss--Hermite quadrature with 30 nodes (Appendix~\ref{app:glmm-fitting}).
To check that 30 nodes were enough, both curvature fits were repeated with 60 nodes, and the coefficients and the standard deviation of the participant intercepts were required to change by less than 0.001.

\clarify{how segments with a missing curvature or bifurcation angle are handled (imputation for the test, a fixed rule in the cross-validation) is not yet decided.}

\subsection{Test of the Geometric Features}
\label{sec:method-test}

Whether the geometric features are associated with incident coronary plaque beyond the established risk factors was assessed with two prespecified likelihood ratio tests, one per model, each of the base against the extended fit (Eq.~\eqref{eq:tech-lr}).
In the curvature model, the statistic was referred to a $\chi^2$ distribution with $q = 1$ degree of freedom.
In the bifurcation model, the bifurcation angle and the daughter asymmetry were tested as a set, so the statistic was referred to a $\chi^2$ distribution with $q = 2$ degrees of freedom.
These are the only two tests in the thesis.
\clarify{significance level, and whether it is adjusted for the two tests.}

From the extended fits, the hazard ratio of each geometric feature was reported per standard deviation with a 95\% confidence interval, computed from the standard error of its coefficient.
These intervals describe the direction and the size of each association and were not used as further tests.
All hazard ratios are adjusted for the other covariates of their model.
In the curvature model, they are also conditional on the participant intercept, as set out in Section~\ref{sec:tech-random-intercept}, so the hazard ratio of curvature was read as an adjusted association among segments with the same frailty.
How much the participants differed beyond their covariates was summarised by the median hazard ratio, computed from the estimated variance of the participant intercepts (Eq.~\eqref{eq:tech-mhr}).
In the bifurcation model, each participant contributes one outcome, so its hazard ratios compare participants and hold no participant intercept.

The reference distribution of the test in the curvature model assumes that the participant intercept captures the dependence between the nine segments (Section~\ref{sec:tech-glmm-likelihood}).
Segments of the same artery may be more alike than segments of different arteries, which a second random effect per artery would allow for (Appendix~\ref{app:glmm-variance}).
\clarify{whether an artery-level random effect is fitted, or the single participant intercept is kept and discussed as a limitation, is not yet decided.}

\subsection{Added Predictive Value}
\label{sec:method-added-value}

As a secondary analysis, it was estimated how much the geometric features add to the prediction of who develops plaque.
In the curvature model, the risk of a participant is the risk of developing plaque in at least one of the nine segments.
Let $A_i = 0$ denote that no segment of participant $i$ developed plaque, and let $H_{ij} = \exp(\eta_{ij})$ be the cumulative hazard of segment $j$ before the participant intercept is applied.
Given $b_i$, the segments were taken as independent, so the probability that none of them developed plaque is the product of their probabilities of no plaque.
Under the cloglog link, this product becomes a sum, and the segment hazards add to one participant hazard,
\begin{equation}
  P(A_i = 0 \mid b_i) = \exp\!\left(-e^{b_i} H_i\right),
  \qquad
  H_i = \sum_{j=1}^{9} H_{ij} .
  \label{eq:method-agg-conditional}
\end{equation}
The participant intercept of a held-out participant is unknown, so this probability was averaged over the distribution of the intercepts (Eq.~\eqref{eq:tech-agg-marginal}), as in Figure~\ref{fig:tech-glmm-marginal}.
The resulting risk was compared with the participant outcome defined in Section~\ref{sec:method-population}.
The bifurcation model gives the risk of plaque in the proximal LAD or the proximal LCx directly (Eq.~\eqref{eq:tech-cloglog}), and this risk was compared with that outcome.
Each risk was computed from the base and from the extended fit of its model.

The risks were summarised by four kinds of measure.
The C-statistic is the probability that a participant who developed plaque was given a higher risk than one who did not, so it describes how well the risks separate the two groups.
Calibration describes whether the risks agree in size with what was observed.
It was summarised by the calibration slope, which is 1 when high and low risks are neither too extreme nor too close together.
The calibration intercept is 0 when the average risk matches the observed proportion.
The Brier score is the mean squared difference between the risk and the outcome.
The log loss is the negative mean log-likelihood of the outcomes under the risks, and a lower value is better for both.
For each measure and model, the difference between the extended and the base fit was reported as an estimate and was not tested.

The two sets of risks were also shown as back-to-back histograms.
If the features add information, the risks of the extended fit spread further from the average risk than those of the base fit.

A small difference was expected in advance for the curvature model.
A segment is only one of the nine terms of the participant hazard (Eq.~\eqref{eq:method-agg-conditional}), so curvature that is clearly associated with plaque in a segment may change the participant risk only slightly.
The added predictive value is therefore reported to describe the size of the associations, while the associations themselves are judged by the likelihood ratio tests.

\subsection{Cross-Validation}
\label{sec:method-cv}

A model fits the data it was estimated on better than it fits new data, so the participant risks behind the measures above were obtained by cross-validation.
The participants were divided into 10 folds, with all nine segments of a participant in the same fold, so that no participant contributed to both the fitting and the evaluation of a model.
The folds were stratified on the participant outcome, which gave each fold about the same proportion of participants who developed plaque.
Each fold was then held out in turn, all four fits were made on the other nine folds, and the participant risks of the held-out fold were computed from them.
Everything estimated from the data was estimated on the training participants alone, namely the means and standard deviations used for standardisation, the coefficients and the variance of the participant intercepts.
The base and the extended fit always used the same folds, so their difference is a paired comparison.

A single division into folds is arbitrary, so the procedure was repeated 50 times with recorded random seeds.
Each measure was computed once per repeat, on the held-out risks of all participants, and reported as the mean over the 50 repeats with the 2.5th and 97.5th percentiles.
These percentiles show how much a measure varies with the division into folds.
\clarify{intervals that reflect the sampling of participants (a bootstrap of the whole procedure) are not yet decided.}

\subsection{Secondary and Sensitivity Analyses}
\label{sec:method-sensitivity}

Beyond the main tests, where in the coronary tree plaque develops was read descriptively from the extended curvature model, outside the likelihood ratio tests.
This reading used the segment intercepts and an added term that lets the hazard ratio of curvature differ between the LAD, the LCx and the RCA.

Three sensitivity analyses examined whether the results depended on choices made in the main analysis.
First, the outcome was restricted to plaque with mild or worse stenosis (Section~\ref{sec:method-population}).
Subsequently, the six risk factors were replaced by the published SCORE2 risk, recalibrated to incident coronary plaque as a single term (Appendix~\ref{app:score2-recal}).
This shows whether the hazard ratios of the geometric features depend on how the risk factors entered.
Finally, diabetes, body mass index and height were added to both base models, to examine whether the geometric features stood in for body size.
